\documentclass[10pt,a4paper]{amsart}
\usepackage[margin=19mm]{geometry}
\usepackage{amsmath,amssymb,mathtools}
\usepackage[scr=rsfs,cal=euler]{mathalfa}
\usepackage{xfrac}
\usepackage[unicode,hidelinks,bookmarks=false]{hyperref}
\newcommand{\ie}{i.e.\ }
\newcommand{\cf}{cf.\ }
\newcommand{\resp}{respectively\ }
\newcommand{\Subsection}{\subsection}
\providecommand{\nobreakdash}{\nobreak\hbox{-}}
\setlength{\emergencystretch}{2em}
\multlinegap0pt
\usepackage{mathtools}
\usepackage{amssymb}
\usepackage{xcolor}
\usepackage[normalem]{ulem}
\usepackage{cancel}
\newtheorem{theorem}{Theorem}
\title{Ramification Subgroups of Knot Groups and their Profinite and Cohomological Structure}
\author{Marina Palaisti, Federico W. Pasini}
\date{Source: arXiv:2605.20365v1 (CC BY 4.0)}
\begin{document}
\maketitle

\noindent\emph{Source and licence.} Ramification Subgroups of Knot Groups and their Profinite and Cohomological Structure. Version arXiv:2605.20365v1 (CC BY 4.0), distributed by arXiv under the Creative Commons Attribution 4.0 International licence, http://creativecommons.org/licenses/by/4.0/, as recorded in the arXiv licence metadata for this exact version and shown on the abstract page https://arxiv.org/abs/2605.20365v1. Full licence notice: \texttt{licenses/arxiv-2605.20365-CC-BY-4.0.txt}.\par\smallskip

\noindent\textit{Editorial scope.} Counted unit: the article's theorem thm:profinite-inertia-rigidity (source0.tex lines 285--293). The article's complete original proof of it is delivered with it (source0.tex lines 295--314, one \texttt{proof} environment of 20 lines). No mathematics is written, repaired or re-derived: the statement and the proof are the article's own lines, in the article's own notation, and no retained line is substituted or re-flowed.  No further source-local context is needed: every cross-reference of every retained line resolves inside the two delivered spans.  Nothing was omitted from any retained span, so the package carries no omission record.  No retained line contains a citation, so the wrapper carries no bibliography and no bibliography entry.  No retained line contains a figure environment, an image-inclusion command or drawing code, so no image is delivered and the package declares no asset dependency.  Editorial formatting only, in addition: an AMS \texttt{amsart} preamble that restates the article's own declarations and macros used by the retained lines, namely the environments \texttt{theorem} on separate counters, together with the standard packages those lines need.  The retained environments print in this document as Theorem 1 in order of appearance, whereas the article prints the same items with the numbering of its own full document.  The complete native source of the article as distributed in the arXiv e-print for this exact version is bundled unchanged as \texttt{sources/200985/source0.tex}.
	\begin{theorem}[Meridian-preserving profinite isomorphisms preserve closed inertia]\label{thm:profinite-inertia-rigidity}
		Assume there is a topological group isomorphism
		\[\Phi:\widehat G \xrightarrow{\ \cong\ } \widehat G',\] such that $\Phi$ sends the conjugacy class of the profinite meridian subgroup $\overline{\langle m\rangle}^{\,\widehat G}$ to the conjugacy class of $\overline{\langle m'\rangle}^{\,\widehat G'}$; equivalently, there exists $y\in \widehat G'$ with
		\[ \Phi\big(\overline{\langle m\rangle}^{\,\widehat G}\big)= y\,\overline{\langle m'\rangle}^{\,\widehat G'}\,y^{-1}.\]
		Let $\widehat U\le \widehat G$ be any open subgroup and set $\widehat U':=\Phi(\widehat U)$.
		Then for every $x\in \widehat G$ one has
		\[ \Phi\big(\widehat I_x(\widehat U)\big)=\widehat I_{\Phi(x)y}(\widehat U')\]
		(where inertia on the right is computed using $m'$). In particular, $\Phi$ induces a bijection between the families of closed inertia subgroups of $\widehat U$ and of $\widehat U'$.
	\end{theorem}

	\begin{proof}
		Fix an open subgroup $\widehat U\le \widehat G$, let $\widehat U'=\Phi(\widehat U)$, and fix $x\in \widehat G$. By definition,
		\[ \widehat I_x(\widehat U)=\widehat U\cap x\,\overline{\langle m\rangle}\,x^{-1}.\]
		Apply $\Phi$ to both sides. Because $\Phi$ is an isomorphism of topological groups, it preserves intersections and conjugation:
		\[ \Phi(A\cap B)=\Phi(A)\cap \Phi(B),\qquad \Phi(xAx^{-1})=\Phi(x)\Phi(A)\Phi(x)^{-1}.\]
		Hence
		\begin{align*}
			\Phi\big(\widehat I_x(\widehat U)\big) &= \Phi(\widehat U)\cap \Phi\big(x\,\overline{\langle m\rangle}\,x^{-1}\big) \\
			&= \widehat U'\cap \Phi(x)\,\Phi\big(\overline{\langle m\rangle}\big)\,\Phi(x)^{-1}.
		\end{align*}
		
		By hypothesis, $\Phi(\overline{\langle m\rangle})=y\,\overline{\langle m'\rangle}\,y^{-1}$ for some $y\in \widehat G'$. (Here $y$ is determined only up to right--multiplication by an element of $\overline{\langle m'\rangle}$, but this ambiguity does not affect the inertia subgroups.) Substituting in the previous equation:
		\begin{align*}
			\Phi\big(\widehat I_x(\widehat U)\big) &= \widehat U'\cap \Phi(x)\,y\,\overline{\langle m'\rangle}\,y^{-1}\,\Phi(x)^{-1} \\
			&= \widehat U'\cap (\Phi(x)y)\,\overline{\langle m'\rangle}\,(\Phi(x)y)^{-1},
		\end{align*}
		which is exactly $\widehat I_{\Phi(x)y}(\widehat U')$ by definition (with meridian $m'$). This proves the claimed equality.
		
		Finally, since $x$ ranges over all of $\widehat G$, so does $\Phi(x)y$ range over all of $\widehat G'$, hence the correspondence is bijective at the level of families.
	\end{proof}

\end{document}
